\section{刘先明的路径：从 CV/ML 到 Robotaxi，再到主机厂}

本章先看刘先明为什么会走到小鹏。他的路径包括 UIUC 博士、Facebook Connectivity Lab、地平线北美、Cruise Robotaxi，再到小鹏。贯穿这条线的不是“跳槽履历”，而是 mission driven 的技术选择：从卫星图像做救灾和世界人口分布，到自动驾驶减少司机疲劳和交通风险，再到主机厂里做可控数据闭环。

\begin{figure}[H]
\centering
\includegraphics[width=0.96\textwidth]{liuxianming-trajectory.png}
\caption{刘先明的 Physical AI 路径：CV/ML、Facebook、Cruise、小鹏和 Physical AI。自制概念图，依据 00:02:16--00:19:00 对谈内容整理。}
\end{figure}

\begin{knowledgebox}{读图：这条路径的主线是真实世界问题}
从卫星图像救灾，到 Robotaxi，再到小鹏量产车队，刘先明一直在处理“AI 如何进入真实世界”。这解释了他为什么不把自动驾驶只看作车企功能，而看作 Physical AI 的一部分。
\end{knowledgebox}

\subsection{Cruise 学到的两件事}

刘先明在 Cruise 学到的两件事影响了他后续路线。第一是极致简化和大规模 Infra。Robotaxi 要从 demo 变成产品，需要 Data Infrastructure、Training Infrastructure、问题分析链路和软件发布体系。第二是 Continuous Learning Machine，即用数据迭代持续解决问题。所谓“躺在夏威夷等金币掉下来”只是玩笑，背后是数据闭环可以让系统持续变好。

\begin{importantbox}{Cruise 经验的迁移}
从 Robotaxi 到主机厂，关键不是换一个场景，而是把“数据回流 + Infra + 快速迭代”迁到更大车队和更可控数据链路中。
\end{importantbox}

\subsection{为什么去主机厂}

如果自动驾驶和 Physical AI 的未来依赖大规模数据、可控数据链路和快速反馈，主机厂就有天然优势。Robotaxi 公司能控制车队，但规模有限；主机厂有量产车辆、真实用户、不同城市道路和完整车端数据链路。刘先明选择小鹏，是因为他和何小鹏都在问同一个问题：下一代技术路线怎样远远甩开当前对手？

\begin{knowledgebox}{主机厂相对 Robotaxi 公司的不同优势}
\begin{tabularx}{\textwidth}{p{0.22\textwidth}p{0.34\textwidth}X}
\toprule
维度 & Robotaxi 公司 & 主机厂 \\
\midrule
车队控制 & 控制运营车队，数据质量高但规模较慢 & 量产车规模大，覆盖城市和场景更广。 \\
商业目标 & 直接提供出行服务 & 卖车、智驾订阅、Robotaxi 和品牌体验并行。 \\
数据链路 & 运营闭环强 & 用户车队回流和量产问题更复杂。 \\
风险约束 & 服务区和运营范围可控 & 面向普通用户，安全和质量责任更重。 \\
\bottomrule
\end{tabularx}
\end{knowledgebox}

\subsection{真实世界问题带来的压力}

刘先明比较了 Facebook、Cruise 和小鹏的工作压力。Facebook 的研究问题更多是技术和资源问题；Cruise 的车在路上，事故、故障和安全都是真实压力；到小鹏之后节奏更快，因为量产车企要同时处理 AI、产品、业务、质量、硬件和用户责任。这种压力解释了为什么 Physical AI 不是纯 research，而是要在真实系统中不断做取舍。

\begin{warningbox}{真实世界 AI 的压力来源}
数字世界模型出错，通常可以回滚、重试或隐藏失败；车在路上出错，关系到安全、责任、品牌和监管。因此物理 AI 的迭代必须比纯软件更谨慎，也更依赖工程闭环。
\end{warningbox}

\subsection{从旧金山到广州：场景改变问题结构}

刘先明提到，在旧金山做自动驾驶时，流浪者、垃圾、宠物、鸟等 corner case 很密集，让人觉得问题很难；而在广州，某些曾经难以处理的问题不是同一个量级。这个观察很重要：自动驾驶不是抽象算法题，城市道路结构、交通参与者、基础设施和人群行为都会改变问题难度。

\begin{importantbox}{场景不是背景，而是模型的一部分}
同一套算法在不同城市会面对不同 corner case。场景选择会影响数据分布、模型训练、上线节奏和用户体验，因此 Robotaxi 和量产智驾都不能只谈模型，不谈城市和运营环境。
\end{importantbox}

\subsection{本章小结}

刘先明的路径说明，小鹏这次换帅不只是组织任命，而是把自动驾驶经验、Robotaxi 数据闭环和主机厂量产资源重新组合。

